\item \textbf{Pauta de evaluación de la Tarea 1}.

\begin{tcolorbox}[colback=blue!2!white,colframe=blue!55!black,title={Criterios generales}]
\begin{itemize}
    \item La evaluación prioriza la coherencia lógica del desarrollo por sobre una única forma de resolver.
    \item Se asigna puntaje parcial a soluciones correctamente encaminadas.
    \item Se permite cualquier estrategia válida que cumpla con los requerimientos del enunciado.
\end{itemize}
\end{tcolorbox}

\begin{solution}

\subsubsection*{Ejercicio 1: Clasificación de número entero (Total: 5 puntos)}

\paragraph{1. Comprensión del problema (1 punto)}
\begin{itemize}
    \item 1.0: Identifica correctamente que debe clasificar par/impar usando módulo.
    \item 0.5: Comprensión parcial.
    \item 0.0: No comprende el objetivo.
\end{itemize}

\paragraph{2. Uso del operador módulo (\%) (1 punto)}
\begin{itemize}
    \item 1.0: Usa correctamente \texttt{numero \% 2}.
    \item 0.5: Uso incorrecto pero con intención clara.
    \item 0.0: No utiliza módulo.
\end{itemize}

\paragraph{3. Uso de estructuras condicionales (1.5 puntos)}
\begin{itemize}
    \item 1.5: Uso correcto de \texttt{if/else}.
    \item 1.0: Condicional presente pero mal implementado.
    \item 0.5: Lógica incompleta.
    \item 0.0: No usa condicionales.
\end{itemize}

\paragraph{4. Implementación en Python (1 punto)}
\begin{itemize}
    \item 1.0: Código correcto y ejecutable.
    \item 0.5: Errores menores.
    \item 0.0: No ejecuta.
\end{itemize}

\paragraph{5. Salida correcta (0.5 puntos)}
\begin{itemize}
    \item 0.5: Mensaje correcto (par/impar).
    \item 0.25: Mensaje ambiguo.
    \item 0.0: No imprime resultado.
\end{itemize}

\vspace{0.5cm}

\subsubsection*{Ejercicio 2: Sumatoria y productoria con iteradores (Total: 5 puntos)}

\paragraph{1. Comprensión del problema (1 punto)}
\begin{itemize}
    \item 1.0: Identifica correctamente sumatoria + productoria.
    \item 0.5: Comprensión parcial.
    \item 0.0: No comprende el objetivo.
\end{itemize}

\paragraph{2. Uso de iteradores (1.5 puntos)}
\begin{itemize}
    \item 1.5: Implementa correctamente dos iteradores independientes.
    \item 1.0: Usa un solo iterador o con errores.
    \item 0.5: Uso incorrecto.
    \item 0.0: No usa iteradores.
\end{itemize}

\paragraph{3. Manejo de acumuladores (1 punto)}
\begin{itemize}
    \item 1.0: Inicializa correctamente (0 para suma, 1 para producto).
    \item 0.5: Error en inicialización.
    \item 0.0: No usa acumuladores correctamente.
\end{itemize}

\paragraph{4. Implementación del cálculo (1 punto)}
\begin{itemize}
    \item 1.0: Calcula correctamente sumatoria y productoria.
    \item 0.5: Errores parciales.
    \item 0.0: Incorrecto.
\end{itemize}

\paragraph{5. Interpretación del resultado (0.5 puntos)}
\begin{itemize}
    \item 0.5: Reconoce que la productoria domina numéricamente.
    \item 0.25: Observación incompleta.
    \item 0.0: No analiza resultado.
\end{itemize}

\vspace{0.5cm}

\subsubsection*{Bonus: uso de funciones (hasta +1 punto)}

\begin{itemize}
    \item +1.0: Implementa correctamente funciones separadas.
    \item +0.5: Funciones incompletas o con errores.
    \item +0.0: No implementa bonus.
\end{itemize}

\vspace{0.5cm}

\subsubsection*{Nota final}

\begin{itemize}
    \item Puntaje base: 10 puntos.
    \item Puntaje máximo con bonus: 11 puntos.
\end{itemize}

\end{solution}
